\answer
\begin{enumerate}

%% 7章 問1（7.1節・単位検算）
\item[7.1節 問1]
$(V_{\mathrm{in}}-V_{\mathrm{out}})$は電圧だからV，$I_{\mathrm{out}}$は電流だからAで，
右辺の単位は
\begin{equation*}
\mathrm{V \cdot A} = \mathrm{W}
\end{equation*}
となり電力の単位に一致する。
$5$\,V$\to 3.3$\,V・$0.5$\,Aの場合は式\ref{eq7.2}より
\begin{equation*}
P_{\mathrm{loss}} = (5 - 3.3)\times0.5 = 0.85\,\mathrm{W}
\end{equation*}
効率は式\ref{eq7.3}より$\eta = 3.3/5 = 0.66 = 66\,\%$である。
入力と出力の電圧差が小さいので，例題\ref{rei7.1}（$42\,\%$）よりずっと効率がよい。

%% 7章 問2（7.2節）
\item[7.2節 問2]
式\ref{eq7.5}より
\begin{equation*}
R_s = \frac{V_{\mathrm{in}} - V_Z}{I_Z}
= \frac{12 - 5.1}{5\times10^{-3}} = 1380\,\Omega
\end{equation*}
おおよそ$1.4$\,k$\Omega$の標準値を選ぶ。
$R_s$が消費する電力は，$R_s$の電圧降下$(V_{\mathrm{in}}-V_Z)=6.9$\,Vに$I_Z$を掛けて
\begin{equation*}
P_{R_s} = 6.9 \times 5\times10^{-3} \approx 35\,\mathrm{mW}
\end{equation*}
である。

%% 7章 問3（7.3節）
\item[7.3節 問3]
式\ref{eq7.9}より，$R_1 = R_2$のときは
\begin{equation*}
V_{\mathrm{out}} = \left(1 + \frac{10}{10}\right)\times2.5 = 2\times2.5 = 5.0\,\mathrm{V}
\end{equation*}
$R_1$を$30$\,k$\Omega$にすると
\begin{equation*}
V_{\mathrm{out}} = \left(1 + \frac{30}{10}\right)\times2.5 = 4\times2.5 = 10\,\mathrm{V}
\end{equation*}
となる。分圧比だけで出力が決まり，$R_1$を上げるほど出力が高くなることがわかる。

\end{enumerate}
